\documentclass{article}
\usepackage{mathtools} 
\usepackage{fontspec}
\usepackage[UTF8]{ctex}
\usepackage{amsthm}
\usepackage{mdframed}
\usepackage{xcolor}
\usepackage{amssymb}
\usepackage{amsmath}
\usepackage{hyperref}
\usepackage{mathrsfs}
\usepackage{mathdots}



% 定义新的带灰色背景的说明环境 zremark
\newmdtheoremenv[
  backgroundcolor=gray!10,
  % 边框与背景一致，边框线会消失
  linecolor=gray!10
]{zremark}{注释}

% 通用矩阵命令: \flexmatrix{矩阵名}{元素符号}{行数}{列数}
\newcommand{\flexmatrix}[4]{
  \[
  #1 = \begin{pmatrix}
    #2_{11}     & #2_{12}     & \cdots & #2_{1#4}   \\
    #2_{21}     & #2_{22}     & \cdots & #2_{2#4}   \\
    \vdots      & \vdots      & \ddots & \vdots     \\
    #2_{#31}    & #2_{#32}    & \cdots & #2_{#3#4}
  \end{pmatrix}
  \]
}

% 简化版命令（默认矩阵名为A，元素符号为a）: \quickmatrix{行数}{列数}
\newcommand{\quickmatrix}[2]{\flexmatrix{A}{a}{#1}{#2}}


\begin{document}
\title{7.2}
\author{张志聪}
\maketitle

\section*{1}

由于$M_k = M_{k + 1}, M_{k - 1} \subset M_k$，于是由命题2.4的证明过程可知，
J中以$\lambda_0$为特征值且阶$\geq k + 1$的Jordan块的个数为
\begin{align*}
  r(I^k) - r(I^{k + 1})
   & = dim M_{k + 1} - dim M_{k} \\
   & = 0
\end{align*}
其中$I = A - \lambda_0 E$。

同理，J中以$\lambda_0$为特征值且阶为$\geq k$的Jordan块的个数为
\begin{align*}
  r(I^{k - 1}) - r(I^{k})
   & = dim M_{k} - dim M_{k - 1} \\
   & \geq 1
\end{align*}
综上可知，$J$中以$\lambda_0$为特征值且阶为$k$的Jordan块的个数至少为1，
且是最高阶数。

\section*{2}

 (i)由题设$M_k = Ker \mathscr{B}^k$可知，
$\forall \alpha \in M_k$，都有
\begin{align*}
  \mathscr{B}^k \alpha & = 0
\end{align*}
于是，对任意$\forall \beta \in N_k$且$\beta \neq 0$，
由于$V = M_k \oplus N_k$，$\beta \notin M_k$，故
\begin{align*}
  \mathscr{B}^k \beta \neq 0
\end{align*}
即
\begin{align*}
  (\mathscr{A} - \lambda_0 E)^k \beta \neq 0 \\
  (\mathscr{A} - \lambda_0 E)^{k-1}(\mathscr{A} - \lambda_0 E) \beta \neq 0
\end{align*}
所以
\begin{align*}
  (\mathscr{A} - \lambda_0 E) \beta & \neq 0               \\
  \mathscr{A} \beta                 & \neq \lambda_0 \beta
\end{align*}
故$\lambda_0$不是$\mathscr{A}\big|_{N_k}$的特征值。


(ii)通过证明$\mathscr{B}\big|_{N_k}$是双射来证明其可逆。

由$N_k = Im(\mathscr{B}^k)$这个题设直接可得$\mathscr{B}\big|_{N_k}$是满射。

单射，通过反证法证明，假设对存在$\alpha, \beta \in N_k$且$\alpha \neq \beta$使得
\begin{align*}
  \mathscr{B} \alpha           & = \mathscr{B} \beta \\
  \mathscr{B} (\alpha - \beta) & = 0
\end{align*}
于是可得$\alpha - \beta \in Ker \mathscr{B}$，这和$V = M_k \oplus N_k$矛盾，
假设不成立，故$\mathscr{B}\big|_{N_k}$是单射。

综上，$\mathscr{B}\big|_{N_k}$可逆。

\section*{3}

由习题2可知，$\lambda_0$不是$\mathscr{A} \big|_{N_k}$的特征值，
于是$\mathscr{A}\big|_{M_k}$中$\lambda_0$的重数，就是$\mathscr{A}$中
$\lambda_0$的重数。

由教材中的内容可知，$\mathscr{B} = \mathscr{A} - \lambda_0 E$
限制在$M_k$内为幂零线性变换，根据命题2.1，在$M_k$中存在一组基，
使得$\mathscr{A}\big|_{M_k}$在这组基下的矩阵成Jordan标准型$J$，
且以$\lambda_0$为特征值的Jordan块，因此由Jordan块的特点，我们有
\begin{align*}
  |\lambda E - J| = (\lambda - \lambda_0)^{dim M_k}
\end{align*}
至此，命题得证。

\section*{4}

$\mathscr{A}$在某组基下的矩阵是Jordan标准型，于是
$\mathscr{A} - \lambda_1 E$在这组基下的矩阵为
\begin{align*}
  J - \lambda_1 E
\end{align*}
于是
\begin{align*}
  (J - \lambda_1 E)^l = \begin{bmatrix}
                          (J_1 - \lambda_1 E)^l &                       &        &                       & 0 \\
                                                & (J_2 - \lambda_1 E)^l &        &                       &   \\
                                                &                       & \ddots &                       &   \\
                          0                     &                       &        & (J_s - \lambda_1 E)^l &   \\
                        \end{bmatrix}
\end{align*}
因为
\begin{align*}
  (J_i - \lambda_1 E)^l & = \begin{bmatrix}
                              \lambda_i - \lambda_1 E & 1                     &        &                       & 0 \\
                                                      & \lambda_i - \lambda_1 & 1      &                       &   \\
                                                      &                       & \ddots & \ddots                &   \\
                              0                       &                       &        & \lambda_i - \lambda_1 & 1 \\
                            \end{bmatrix}^l
\end{align*}

(i)若$\lambda_i \neq \lambda_1$，则$(J_i - \lambda_1 E)^l$满秩，没有非零解；

(ii)若$\lambda_i = \lambda_1$，则$(J_i - \lambda_1 E)^l$不满秩；

回到本题中的$\alpha$，如果$\alpha = 0$，则$\alpha \in N_k$，结论成立；

$\alpha \neq 0$，设$\alpha$在这组基下的坐标为$X$，由题设
\begin{align*}
  (\mathscr{A} - \lambda_1 E)^l \alpha & = 0 \\
  (J - \lambda_1 E)^l X                & = 0
\end{align*}
于是$(J_i - \lambda_1 E)^l, i \neq 1$是满秩的，没有非零解，
故坐标$X$在所有特征值$\lambda_i \neq \lambda_1$的Jordan块上的坐标分量全部为0.

反证法，假设$\alpha \in M_k$，即
\begin{align*}
  (\mathscr{A} - \lambda_0 E)^k \alpha & = 0
\end{align*}
同理可得，坐标$X$在所有特征值$\lambda_i \neq \lambda_0$的Jordan块上的坐标分量全部为0.

综上，$X = 0$，这与$\alpha \neq 0$矛盾，故假设不成立，又因为$V = M_k \oplus N_k$，
故$\alpha \in N_k$.

\section*{5}

在复数域上，任何方阵都相似于Jordan标准型，即存在可逆矩阵P，使得
\begin{align*}
  P^{-1} A P & = J
\end{align*}
其中
$J = \begin{pmatrix}
    J_1 &        &     \\
        & \ddots &     \\
        &        & J_s
  \end{pmatrix}$，
由题设$A^k = E$可知
\begin{align*}
  J^k = (P^{-1} A P)^k = E
\end{align*}
接下来证明J是对角矩阵，即J的Jordan块都是一阶的。

反证法，假设$J$存在大于1阶的Jordan块$J_i$，令
\begin{align*}
  J_i = \lambda E + N
\end{align*}
其中$N$是第一条上对角线元素都是1，于是由矩阵乘法可知
\begin{align*}
  J_i^k                                                                            & = E_i \\
  (\lambda E + N)^k                                                                & = E_i \\
  \lambda^k E_i + \lambda^{k - 1} E_i N + C_{k}^2 \lambda^{k - 2} E_i N^2 + \cdots & = E_i
\end{align*}
由于$\lambda^k E_i$没有上对角线，$N^n, n = 2, 3, \cdots$的第一条上对角线元素都是0，
所以$J_i^k \neq E_i$，出现矛盾，所以$J$的Jordan块都是一阶的。

综上，命题得证。

\section*{6}
略

\section*{7}

显然，$A$的特征值a具有n重，
由习题3可知，$dim M_k$也是n，
于是由习题1可知，J中以a为特征值的Jordan块是n阶的，
故$J$只有一个Jordan块，且是n阶的，所以，A的Jordan标准型是
\begin{align*}
  J = \begin{bmatrix}
        a & 1 &        &        & 0 \\
          & a & 1      &        &   \\
          &   & \ddots & \ddots &   \\
          &   &        & a      & 1 \\
          &   &        &        & a
      \end{bmatrix}
\end{align*}

\section*{8}

J是幂零矩阵，当$k \geq n$，则$J^k = 0$，
此时，Jordan标准型就是零矩阵。

当$k < n$时，由P79关于幂零矩阵的乘法性质，我们有
$J^k$是上三角矩阵，且主对角线元素都是0，
于是$J^k$的特征多项式
\begin{align*}
  f(\lambda) & = |\lambda E - J^k| \\
             & = \lambda^n
\end{align*}
由第一节习题1可知，$J^k$是幂零矩阵。

接下来，考虑Jordan块的阶数问题。

按命题2.1中的方式，
令$B = J^k - \lambda E = J^k - 0 \cdot E = J^k$，
$M_i, N_i, i = 0, 1, \cdots$的定义与教材中一致。

于是，我们有
\begin{align*}
  B^l = (J^k)^l = J^{lk}
\end{align*}
使用整数的带余除法，我们有
\begin{align*}
  n = q k + r
\end{align*}
于是，当$l \geq q$时，$J^{lk} = 0$，
此时$dim M_l = dim Ker B^l = dim Ker J^{lk} = n$；
$l < q$时，$dim B^l = dim J^{lk} = n - lk$，于是
$dim M_l = n - dim B^l = lk$.

利用命题2.4，我们有如下结论。
\begin{itemize}
  \item $l + 1 \leq q$时，$l$阶的Jordan块的个数为
        \begin{align*}
          2 dim M_l - dim M_{l + 1} - dim M_{l - 1}
           & = 2 l k - (l + 1) k - (l - 1) k \\
           & = 2lk - l k - k - lk + k        \\
           & = 0
        \end{align*}

  \item $l = q$时，$l$阶的Jordan块的个数为
        \begin{align*}
          2 dim M_l - dim M_{l + 1} - dim M_{l - 1}
           & = 2 lk - n - (l - 1) k        \\
           & = 2 qk - (qk + r) - (q - 1) k \\
           & = 2 qk - qk - r - qk + k      \\
           & = k - r
        \end{align*}

  \item $l = q + 1$时，$l$阶的Jordan块的个数为
        \begin{align*}
          2 dim M_l - dim M_{l + 1} - dim M_{l - 1}
           & = 2 n - n - (l - 1) k     \\
           & = 2 n - n - (q + 1 - 1) k \\
           & = 2 n - n - q k           \\
           & = n - qk                  \\
           & = r
        \end{align*}

  \item $l > q + 1$时，$l$阶的Jordan块的个数为
        \begin{align*}
          2 dim M_l - dim M_{l + 1} - dim M_{l - 1}
           & = 2 n - n - n \\
           & = 0
        \end{align*}
\end{itemize}

综上，$J^k$的幂零矩阵中有$k - r$个$q$阶Jordan块，$r$个$q + 1$阶Jordan块。

注：一共有$k$个Jordan块，这与$dim J^k = n - k$，$J^k$的解空间的维数是$n - dim J^k = k$是一致的。
回顾第一节习题10.

\section*{9}

\begin{itemize}
  \item (1) 方法一;

        通过把$A$和$A^T$都看作在某组基下线性变换的矩阵，
        两者的Jordan标准型$J$相同。

        由命题2.2和命题2.4可知，线性变换的Jordan标准型只和
        它的特征值和每个Jordan块的阶有关（即和$dim M_i$相关）。

        我们知道，转置不影响行列式的值，即
        \begin{align*}
          |\lambda E - A|
           & = |(\lambda E - A)^T| \\
           & = |\lambda E - A^T|
        \end{align*}
        所以，两者的特征值相同。

        对特征值$\lambda_0$，令
        \begin{align*}
          B = A - \lambda_0 E
        \end{align*}
        于是
        \begin{align*}
          rank (A - \lambda_0 E)^k = rank (A^T - \lambda_0 E)^k,
          k = 1, 2, \cdots
        \end{align*}
        以上结论通过对k进行归纳容易得到。

        于是，我们有
        \begin{align*}
          dim (Ker B^i)
           & = n - dim (A - \lambda_0 E)^i   \\
           & = n - dim (A^T - \lambda_0 E)^i
        \end{align*}

        综上，命题得证。

  \item (2) 方法二.

        复数矩阵$A$与Jordan标准型相似，即存在可逆矩阵使得
        \begin{align*}
          A = P^{-1} J P
        \end{align*}
        于是
        \begin{align*}
          A^T & = (P^{-1} J P)^T     \\
              & = P^T J^T (P^{-1})^T
        \end{align*}
        令$Q = (P^{-1})^T$，则
        \begin{align*}
          A^T = Q^{-1} J^T Q
        \end{align*}
        所以，$A^T \backsim  J^T$.

        利用反序矩阵的性质，令
        \begin{align*}
          R = \begin{bmatrix}
                0      & 0       & 0      & 1      \\
                0      & 0       & 1      & 0      \\
                \vdots & \iddots & \ddots & \vdots \\
                1      & 0       & 0      & 0
              \end{bmatrix}_{n \times n}
        \end{align*}
        于是，我们有
        \begin{align*}
          J^T = R^{-1} J R
        \end{align*}
        （注：可以通过一个特别的例子来看具体的过程。）
        所以，$J \backsim J^T$。

        综上，$A^T \backsim A$.
\end{itemize}

\section*{10}

\begin{itemize}
  \item (1)

        (i) $M_i$是子空间；

        加法封闭:
        任意$\alpha, \beta \in M_i$，则有存在$m_1, m_2$使得
        \begin{align*}
          (\mathscr{A} - \lambda_i E)^{m_1} \alpha & = 0 \\
          (\mathscr{A} - \lambda_i E)^{m_2} \beta  & = 0
        \end{align*}
        取$m = max\{m_1, m_2\}$，于是
        \begin{align*}
          (\mathscr{A} - \lambda_i E)^{m} (\alpha + \beta)
           & = (\mathscr{A} - \lambda_i E)^{m} \alpha + (\mathscr{A} - \lambda_i E)^{m} \beta                \\
           & = (\mathscr{A} - \lambda_i E)^{m - m_1} (\mathscr{A} - \lambda_i E)^{m_1} \alpha
          + (\mathscr{A} - \lambda_i E)^{m - m_2} (\mathscr{A} - \lambda_i E)^{m_2} \beta                    \\
           & = (\mathscr{A} - \lambda_i E)^{m - m_1} \cdot 0 + (\mathscr{A} - \lambda_i E)^{m - m_2} \cdot 0 \\
           & = 0
        \end{align*}
        故$\alpha + \beta \in M_i$.

        数乘封闭:
        任意$\alpha, \beta \in M_i$，因为$(\mathscr{A} - \lambda_i E)^{m}$
        是线性变换，所以
        \begin{align*}
          (\mathscr{A} - \lambda_i E)^{m} (k \alpha)
           & = k (\mathscr{A} - \lambda_i E)^{m} \alpha \\
           & = k \cdot 0                                \\
           & = 0
        \end{align*}
        故$k \alpha \in M_i$.

        (ii) $M_i$是不变子空间。

        对任意$\alpha \in M_i$，有
        \begin{align*}
          (\mathscr{A} - \lambda_i E)^{m} \alpha = 0
        \end{align*}
        另外
        \begin{align*}
          (\mathscr{A} - \lambda_i E)^m (\mathscr{A} \alpha)
           & = \mathscr{A} (\mathscr{A} - \lambda_i E)^m \alpha \\
           & = \mathscr{A} \cdot 0                              \\
           & = 0
        \end{align*}
        注：第一个等式可以通过对m进行归纳完成证明。

        故$\mathscr{A} \alpha \in M_i$，所以$M_i$是不变子空间。
  \item (2)

        因为$\mathscr{A}$的特征值都在数域K上，所以存在一组基，使得$\mathscr{A}$在这组基下的矩阵为Jordan标准型$J$.
        在这组基下：
        \begin{align*}
          V = V_1 \oplus V_2 \oplus \cdots \oplus V_k
        \end{align*}
        其中$V_i:$由所有以$\lambda_i$为特征值的Jordan块对应的基向量张成。

        命题转换为证明$V_i = M_i$.

        (i)在$V_i$上，在它的基向量下，对应的矩阵
        也是Jordan标准型，且每个Jordan块都有
        \begin{align*}
          J_{\lambda_i} = \lambda_i E + N
        \end{align*}
        其中N是幂零矩阵。所以
        \begin{align*}
          (\mathscr{A} - \lambda_i \mathscr{E}) \big|_{V_i}
        \end{align*}
        的矩阵为$N$，而$N^r = 0$（某个整数$r$）.
        于是对任意$\alpha \in V_i$:
        \begin{align*}
          (\mathscr{A} - \lambda_i \mathscr{E})^r \alpha = 0
        \end{align*}
        故$V_i \subseteq M_i$.

        (ii)反过来，任取$\alpha \in M_i$，把$V$拆分为
        \begin{align*}
          V = V_i \oplus \left( \bigoplus \limits_{j \neq i} V_j \right)
        \end{align*}
        于是$\alpha$可表示为
        \begin{align*}
          \alpha = \alpha_i + \sum\limits_{j \neq i} \alpha_j
        \end{align*}
        其中$\alpha_i \in V_i$，$\alpha_j \in V_j$.

        因为$\alpha \in M_i$，存在$m$使得
        \begin{align*}
          (\mathscr{A} - \lambda_i \mathscr{E})^m \alpha = 0
        \end{align*}
        即
        \begin{align*}
          (\mathscr{A} - \lambda_i \mathscr{E})^m (\alpha_i + \sum\limits_{j \neq i} \alpha_j) & = 0 \\
          (\mathscr{A} - \lambda_i \mathscr{E})^m \alpha_i
          + \sum\limits_{j \neq i} (\mathscr{A} - \lambda_i \mathscr{E})^m \alpha_j            & = 0
        \end{align*}
        因为$V_i, V_j (j \neq i)$都是不变子空间，所以
        \begin{align*}
          (\mathscr{A} - \lambda_i \mathscr{E})^m \alpha_i \in V_i \\
          (\mathscr{A} - \lambda_i \mathscr{E})^m \alpha_j \in V_j, j \neq i
        \end{align*}
        由直和0向量表示的唯一性可知
        \begin{align*}
          (\mathscr{A} - \lambda_i \mathscr{E})^m \alpha_i = 0 \\
          (\mathscr{A} - \lambda_i \mathscr{E})^m \alpha_j = 0, j \neq i
        \end{align*}
        又因为，在$V_j$上，$\mathscr{A} - \lambda_i \mathscr{E}$
        的矩阵为$J_j - \lambda_i E$是满秩的，故$\alpha_j = 0$，
        所以
        \begin{align*}
          \alpha = \alpha_i \in V_i
        \end{align*}
        于是可得 $M_i \subseteq V_i$.

        综上可得，$V_i = M_i$.
  \item (3)

        (i) $\mathscr{A}\big|_{M_i}$的Jordan标准型，就是保留以$\lambda_i$为特征值Jordan块；

        (ii) 由(2)的证明过程可知$V_i = M_i$，为了讨论方便不妨设
        $\bigoplus \limits_{j \neq i} V_j$的基向量为
        $\epsilon_1, \epsilon_2, \cdots, \epsilon_t$，
        于是可得
        \begin{align*}
          \overline{\epsilon_1}, \overline{\epsilon_2}, \cdots, \overline{\epsilon_t}
        \end{align*}
        是$V/M_i$的一组基，在这组基下，$\mathscr{A}$在$V/M_i$内的诱导变换的Jordan标准型
        就是$J$中去掉以$\lambda_i$为特征值Jordan块。
\end{itemize}

\end{document}
